\documentclass[11pt]{article}
\usepackage{fontspec}
\usepackage[french]{babel}
\usepackage[leqno]{amsmath}
\usepackage{amssymb}
\usepackage{amsthm}
\usepackage{mathrsfs}
\usepackage{tikz-cd}
\usepackage{enumitem}

\newcommand{\printedpage}[1]{\marginpar{\footnotesize p.~#1}}
\newcommand{\msp}[1]{\marginpar{\footnotesize [#1]}}
\newcommand{\Hot}{\mathrm{Hot}}
\newcommand{\Ch}{\underline{\mathrm{Ch}}}
\newcommand{\Ob}{\mathrm{Ob}\,}
\newcommand{\Cat}{\mathrm{Cat}}

\begin{document}

\printedpage{4}
% head: A. Grothendieck : Dérivateurs, Ch. XI, Édition des Universités de Montpellier II et Paris VII
% foot: Version 28 mai 2002

\textbf{Proposition 1.} \textit{Cette notion est stable par changement de base (trivial), par composition (formel sur a)). Pour que $f$ soit une \textup{Hot}-fibration, il faut et il suffit que $f$ \textup{\texttt{[plutôt}} $f^{o}$\textup{\texttt{]}} le soit.}

\textbf{Proposition 2.} \textit{Soit $f : X \longrightarrow S$ une \textup{Hot}-fibration. Alors $s \longmapsto \mathrm{hot}(X_s)$ est un système local sur $S$, à valeurs dans \textup{Hot}. Donc les $X_s$ pour $s$ dans une même composante connexe de $S$ sont \textup{Hot}-isomorphes entre eux (mais pas canoniquement, si cette composante n'est pas \textup{1}-connexe).}

% (a vertical change bar is printed in the left margin beside Théorème 2)
\textbf{Théorème 2.} {\itshape Soit
\begin{tikzcd}[ampersand replacement=\&, baseline=(\tikzcdmatrixname-1-1.base), column sep=small, row sep=small]
X' \arrow[rr, "f"] \arrow[dr, "p'"'] \& \& X \arrow[dl, "p"] \\
\& S \&
\end{tikzcd}
un diagramme commutatif dans \textup{Cat}, avec $p$, $p'$ des \textup{Hot}-fibrations, et soit $S_0 \subseteq S$ une partie de $S$ telle que $S_0 \longrightarrow \pi_0(S)$ soit surjectif. Pour que $f$ soit une \textup{Hot}-équivalence, il faut et il suffit que les $f_s : X_s \longrightarrow Y_s$ $(s \in S_0)$ le soient.} % [sic] f_s : X_s -> Y_s -> f_s : X'_s -> X_s, the diagram has X' and X (cf. f_s : X'_s -> X_s p. 6)

a) La suffisance se prouve comme tantôt, avec le critère de Artin-Mazur.

\msp{page 132}

Notons que l'hypothèse sur $S_0$ assure que si $F$, $F'$ sont deux faisceaux d'ensembles localement constants sur $S$, et $u : F \longrightarrow F'$ un homomorphisme tel que $u_s : F_s \xrightarrow{\ \sim\ } F'_s$ pour $s \in S$, alors $u$ est \texttt{[un]} isomorphisme. % [sic] s \in S -> s \in S_0, otherwise the hypothesis on S_0 plays no role

1$^{o}$) Soit $F$ un faisceau d'ensembles localement constant sur $X$, $F'$ son image inverse sur $X'$, il faut prouver
\[
H^0(X,F) \ \xrightarrow{\ \sim\ }\ H^0(X',F').
\]
Or
\[
H^0(X,F) \simeq H^0(Y, \underbrace{f_*(F)}_{\substack{\text{localement}\\ \text{constant sur } S}}) \ \longrightarrow\ H^0(Y, f'_*(F')) \simeq H^0(X',F'),
\] % [sic] H^0(Y, f_*(F)), H^0(Y, f'_*(F')) -> H^0(S, p_*(F)), H^0(S, p'_*(F')), the brace says "localement constant sur S" and (*) uses p_*, p'_*
il faut prouver que
\[
p_*(F) \ \xrightarrow{\ \sim\ } p'_*(F') \tag*{$(*)$}
\]
est \texttt{[un]} isomorphisme. Or par hypothèse sur $p$, $p'$, les deux faisceaux sont localement constants et se calculent fibre par fibre, d'où résulte que $(*)$ est un isomorphisme pour les fibres en $s \in S_0$, donc un isomorphisme partout. OK.

2$^{o}$) Il en résulte déjà, si $F$ est un faisceau en \textit{groupes,} que le foncteur
\[
(F\text{-torseurs sur } X) \ \longrightarrow\ (F\text{-torseurs sur } X')
\]
est pleinement fidèle, donc $H^1(X,F) \longrightarrow H^1(X',F')$ \texttt{[est]} injectif; il reste à prouver que ce foncteur est essentiellement surjectif. À cause de la pleine fidélité, c'est \textit{local sur $S$} (\textbf{NB} On voit directement que l'hypothèse fait sur les $s \in S_0$ implique la même hypothèse pour tout $s \in S$, en utilisant le fait que $p$, $p'$ sont des fibrations) (\footnote{Les hypothèses faites sont bel et bien locales sur $S$ (du moins si $S_0 = S$).}). Donc on est ramené au cas où $S = S/s_0$ a un objet final $s_0$.

\printedpage{5}
% head: A. Grothendieck : Dérivateurs, Ch. XI, Édition des Universités de Montpellier II et Paris VII
% foot: Version 28 mai 2002

\msp{page 133}

On a alors
\[
\begin{aligned}
H^1(X,F) &\ \simeq\ H^0(S, R^1p_*(F)) \\
H^1(X',F') &\ \simeq\ H^0(S, R^1p'_*(F')),
\end{aligned}
\]
et on voit encore que
\[
R^1p_*(F) \ \longrightarrow\ R^1p'_*(F')
\]
est \texttt{[un]} isomorphisme (car il se calcule fibre par fibre, et on utilise l'hypothèse sur les $f_s$), donc induit un isomorphisme
\[
H^1(X,F) \ \longrightarrow\ H^1(X',F').
\]

3$^{o}$) Si $F$ est même abélien, on doit prouver
\[
H^i(X,F) \ \xrightarrow{\ \sim\ }\ H^i(X',F')
\]
pour tout $i$, et cela résulte encore de la suite spectrale de Leray pour $(p,F)$ et pour $(p',F')$.

b) Inversement, supposons que $f$ soit une Hot-équivalence, et prouvons que les $f_s : X'_s \longrightarrow X_s$ sont des Hot-équivalences. Je vais utiliser le

% (a vertical change bar is printed in the left margin beside the Lemme)
\textbf{Lemme.} \textit{Soit $S \in \Ob\Cat$, $S$ \textup{0}-connexe. Alors il existe un $\overline{S} \in \Ob\Cat$, $\overline{S}$ \textup{Hot}-asphérique, et $\overline{S} \longrightarrow S$ une \textup{Hot}-fibration.}

Voyons comment ceci implique le théorème 2 (partie réciproque). On peut supposer $S$ 0-connexe, et on fait le changement de base par $\overline{S} \longrightarrow S$

\msp{page 134}

comme dans le lemme :
\[
\begin{tikzcd}[column sep=small, row sep=small]
X' \arrow[rr, "f"] \arrow[dr, "p'"'] & & X \arrow[dl, "p"] & & & & \\
& S & & & \overline{X'} \arrow[rr, "\overline{f}"] \arrow[dr, "\overline{p'}"'] & & \overline{X} \arrow[dl, "\overline{p}"] \\
& & & & & \overline{S} \arrow[ullll, "\varphi"] &
\end{tikzcd}
\quad .
\]

Alors $\overline{p'}$ et $\overline{p}$ sont des Hot-fibrations, et comme $\varphi$ est une Hot-fibration, le changement de base par $\varphi$ transforme la Hot-équivalence $f$ en une Hot-équivalence $\overline{f}$. Considérons alors $\overline{s} \in \Ob\overline{S}$, et le diagramme
\[
\begin{tikzcd}
\overline{X'} \arrow[r, "\overline{f}"] & \overline{X} \\
\overline{X'}_{\overline{s}} \arrow[u, hook, "i'"] \arrow[r, "\overline{f}_{\overline{s}}"] & \overline{X}_{\overline{s}} \arrow[u, hook, "i"']
\end{tikzcd}
\qquad .
\]

\printedpage{6}
% head: A. Grothendieck : Dérivateurs, Ch. XI, Édition des Universités de Montpellier II et Paris VII
% foot: Version 28 mai 2002

Comme $\{\overline{s}\} \longrightarrow \overline{S}$ est une Hot-équivalence, la flèche $i$ qui s'en déduit par le changement de base $\overline{p} : \overline{X} \longrightarrow \overline{S}$ l'est aussi, puisque $\overline{p}$ est une Hot-fibration, et de même pour $i'$. Ainsi dans le carré commutatif précédent, $\overline{f}$, $i$ et $i'$ sont des Hot-équivalences, donc aussi $\overline{f}_{\overline{s}}$. Mais si $s = \varphi(\overline{s})$, $\overline{f}_{\overline{s}}$ s'identifie à $f_s : X'_s \longrightarrow X_s$. D'ailleurs, il est immédiat que $\varphi$, étant une Hot-fibration

\msp{page 135}

à source non-vide, et à but connexe, est surjective. (\textbf{NB} toutes les fibres d'une même composante connexe de $S$ sont Hot-équivalentes, donc si l'une est non-vide, toutes le sont.) Cela achève la démonstration, modulo le lemme.

Il reste à prouver celui-ci. Comme $S = \emptyset$ \texttt{[plutôt} $S \neq \emptyset$\texttt{]}, je peux choisir une ‘origine’
\[
s_0 \in S,
\]
et regarder la catégorie
\[
\Ch_{\infty}(S; s_0, -)
\]
des chemins infinis dans $S$, d'origine $s_0$, fibre de $\Ch_{\infty}(S)$ par le foncteur-source
\[
\sigma : \Ch_{\infty}(S) \ \longrightarrow\ S.
\]
On sait que les fibres de $\sigma$ sont Hot-asphériques (VII), et la catégorie construite ici n'est autre que (avec les notations de VII, p.~152) que la catégorie $\overline{S'}$ qui entre dans la factorisation canonique de
\[
f : \{s_0\} = S' \ \longrightarrow\ S
\]
en
\[
\{s_0\} = S' \ \underbrace{\longrightarrow}_{\text{Hot-équivalence}}\ \overline{S'} \ \xrightarrow{\ \overline{f}\ }\ S,
\]
où ici $\overline{f}$ s'interprète simplement comme le foncteur but (i.e.\ le foncteur induit sur

\msp{page 136}

$\Ch_{\infty}(S; s_0, -)$ par le foncteur but
\[
\beta : \Ch_{\infty}(S) \ \longrightarrow\ S \ \ ).
\]
On a prouvé dans loc.\ cit.\ que $\overline{f}$ est une Hot-fibration (‘propre’ et ‘lisse’), mais dans un contexte qui semblait faire appel tacitement à la théorie du dérivateur HOT. Donc pour mettre les résultats ‘topologiques’ de loc.\ cit.\ sur une base solide, il faut développer les notions pertinentes de foncteur propre, foncteur lisse, sans référence à une théorie HOT. C'est ce que j'avais d'ailleurs fait déjà en 1983, du temps de Pursuing Stacks, mais sans rédaction en forme à ce sujet. Il me faut donc à présent revenir sur ces notions, de façon un peu systématique. D'ailleurs, j'aurais dû mal à construire un seul foncteur Hot-fibrant (sauf les exemples triviaux, tels les isomorphismes, ou les produits $Y \times Z \xrightarrow{\mathrm{pr}_1} Y$), si je ne disposais d'une théorie des foncteurs lisses et des foncteurs propres (\footnote{J'exagère un peu quand même -- après tout il doit être assez évident que les foncteurs de Thomason-Kan sont des Hot-fibrations (mais il y en a encore un peu trop peu !).}). % [sic] dû mal -> du mal, "avoir du mal à" (partitive article, no accent)

\end{document}
